% ---
% titre: Calcul flash: 9VP - cycle 2 (semaines 7-12) - théorie
% theme: NO
% chapitre: Nombres naturels et décimaux
% sous_chapitre: Calcul réfléchi et mental
% niveau: [9e]
% type: TH
% tags: [c-nombres-decimaux,c-droite-numerique,c-operations,c-addition,c-soustraction,c-criteres-divisibilite,c-decomposition-facteurs-premiers,c-nombres-premiers,c-ppmc, f-individuel,p-flash-calcul]
% difficulte: 1
% source: latex
% ---

\semaineinvisible
\header{Calcul mental: 9VP - Semaine 8 - Décimaux : droite graduée} 

\vfill
%-----------------------------------------------------------------------
\subsection{Placer un nombre décimal}

\begin{regle}{Définition}
Sur une droite graduée, chaque graduation représente un même écart. On repère d'abord entre quels nombres entiers se situe le décimal, puis on affine selon le pas de graduation.
\end{regle}

\vspace{6pt}
\begin{astuce}{Stratégie}
\textit{Exemple :} pour placer $2,4$ entre $2$ et $3$ (graduation par dixièmes), on compte $4$ petits traits après $2$.
\end{astuce}

%-----------------------------------------------------------------------
\vspace{10pt}
\subsection{Comparer deux décimaux}

\begin{regle}{Définition}
On compare d'abord la partie entière. Si elle est égale, on compare les dixièmes, puis les centièmes, etc.
\end{regle}

\vspace{6pt}
\begin{astuce}{Stratégie}
\textit{Exemple :} $3,45 < 3,5$ car $3,45 = 3,450$ et $450 < 500$ millièmes.
\end{astuce}

%-----------------------------------------------------------------------
\vspace{10pt}
\subsection{Encadrer un décimal}

\begin{regle}{Définition}
Encadrer un nombre, c'est trouver les deux multiples les plus proches (d'une unité, d'un dixième, etc.) entre lesquels il se situe.
\end{regle}

\vspace{6pt}
\def\arraystretch{1.2}
\begin{tabularx}{\linewidth}{|X|X|}
\hline
\cellcolor{themeNO!30}\textbf{Nombre} & \cellcolor{themeNO!30}\textbf{Encadrement} \\\hline
$7,3$ (à l'unité) & $7 < 7,3 < 8$ \\\hline
$7,36$ (au dixième) & $7,3 < 7,36 < 7,4$ \\\hline
\end{tabularx}
\def\arraystretch{1}


\newpage
\semaineinvisible
\header{Calcul mental: 9VP - Semaine 9 - Compléments} 

\vfill
%-----------------------------------------------------------------------
\subsection{Le complément}

\begin{regle}{Définition}
Le complément d'un nombre à 10, 100 ou 1000, c'est ce qu'il faut lui ajouter pour atteindre ce total.
\end{regle}

\vspace{6pt}
\begin{astuce}{Stratégie}
\textit{Exemple :} le complément de $37$ à $100$ est $63$, car $37 + 63 = 100$.
\end{astuce}

%-----------------------------------------------------------------------
\vspace{10pt}
\subsection{Stratégie de calcul}

\begin{astuce}{Méthode}
On complète d'abord au dizaine/centaine la plus proche, puis on ajuste.\\
\textit{Exemple :} complément de $68$ à $100$ : $68 \to 70$ (il manque $2$), puis $70 \to 100$ (il manque $30$) $\to 2 + 30 = 32$.
\end{astuce}

%-----------------------------------------------------------------------
\vspace{10pt}
\subsection{Compléments à 1000}

\begin{astuce}{Stratégie}
\textit{Exemple :} complément de $640$ à $1000$ : $640 \to 700$ (il manque $60$), puis $700 \to 1000$ (il manque $300$) $\to 60 + 300 = 360$.
\end{astuce}

%-----------------------------------------------------------------------
\vspace{10pt}
\subsection{Exemples}

\def\arraystretch{1.2}
\begin{tabularx}{\linewidth}{|X|X|}
\hline
\cellcolor{themeNO!30}\textbf{Nombre} & \cellcolor{themeNO!30}\textbf{Complément à 100} \\\hline
$45$ & $55$ \\\hline
$82$ & $18$ \\\hline
$9$ & $91$ \\\hline
\end{tabularx}
\def\arraystretch{1}


\newpage
\semaineinvisible
\header{Calcul mental: 9VP - Semaine 10 - Critères de divisibilité} 

\vfill
%-----------------------------------------------------------------------
\subsection{Qu'est-ce qu'un critère de divisibilité ?}

\begin{regle}{Définition}
Un critère de divisibilité permet de savoir si un nombre est divisible par un autre \textbf{sans faire la division}, en observant simplement ses chiffres.
\end{regle}

%-----------------------------------------------------------------------
\vspace{10pt}
\subsection{Les critères à connaître}

\def\arraystretch{1.3}
\begin{tabularx}{\linewidth}{|c|X|X|}
\hline
\cellcolor{themeNO!30}\textbf{Divisible par} & \cellcolor{themeNO!30}\textbf{Critère} & \cellcolor{themeNO!30}\textbf{Exemple} \\\hline
$2$ & Le dernier chiffre est pair ($0,2,4,6,8$) & $348$ \\\hline
$5$ & Le dernier chiffre est $0$ ou $5$ & $345$ \\\hline
$10$ & Le dernier chiffre est $0$ & $340$ \\\hline
$3$ & La somme des chiffres est divisible par $3$ & $324 \to 3{+}2{+}4=9$ \\\hline
$9$ & La somme des chiffres est divisible par $9$ & $324 \to 3{+}2{+}4=9$ \\\hline
$4$ & Le nombre formé par les 2 derniers chiffres est divisible par $4$ & $512 \to 12$ \\\hline
$25$ & Le nombre formé par les 2 derniers chiffres est divisible par $25$ & $375 \to 75$ \\\hline
$100$ & Les 2 derniers chiffres sont $00$ & $2400$ \\\hline
$6$ & Divisible par $2$ \textbf{et} par $3$ & $324$ \\\hline
$50$ & Divisible par $2$ \textbf{et} par $25$ & $350$ \\\hline
\end{tabularx}
\def\arraystretch{1}

%-----------------------------------------------------------------------
\vspace{10pt}
\subsection{Combiner les critères}

\begin{astuce}{Stratégie}
Pour un critère combiné (comme $6$ ou $50$), on vérifie les deux critères séparément : le nombre doit remplir les deux conditions à la fois.
\end{astuce}


\newpage
\semaineinvisible
\header{Calcul mental: 9VP - Semaine 11 - Décomposition, PGCD \& PPCM} 

\vfill
%-----------------------------------------------------------------------
\subsection{Rappel : décomposition en facteurs premiers}

\begin{regle}{Définition}
Décomposer un nombre en facteurs premiers, c'est l'écrire comme un produit de nombres premiers.
\end{regle}

\vspace{6pt}
\begin{astuce}{Stratégie}
\textit{Exemple :} $60 = 2 \times 2 \times 3 \times 5 = 2^2 \times 3 \times 5$
\end{astuce}

%-----------------------------------------------------------------------
\vspace{10pt}
\subsection{Rappel : PGCD}

\begin{regle}{Définition}
Le \textbf{PGCD} (plus grand commun diviseur) de deux nombres est le plus grand nombre qui les divise tous les deux.
\end{regle}

\vspace{6pt}
\begin{astuce}{Stratégie}
\textit{Exemple :} $\text{PGCD}(12, 18) = 6$
\end{astuce}

%-----------------------------------------------------------------------
\vspace{10pt}
\subsection{Rappel : PPCM}

\begin{regle}{Définition}
Le \textbf{PPCM} (plus petit commun multiple) de deux nombres est le plus petit nombre qui est un multiple des deux.
\end{regle}

\vspace{6pt}
\begin{astuce}{Stratégie}
\textit{Exemple :} $\text{PPCM}(4, 6) = 12$
\end{astuce}